\chapter{Die Kommandozeile}
\label{kap:kommandozeile}

\section{Aufruf}

\begin{shell}
esim <datei.json> [-o <ordner>] [--summary <datei.json>] [-j <n>] [--progress]
\end{shell}

\begin{center}
\begin{tabularx}{\textwidth}{@{}P{3.6cm}L@{}}
  \toprule
  Option & Wirkung \\
  \midrule
  \code{-o <ordner>} & Ordner für die Ergebnisdateien (Standard: das aktuelle Verzeichnis); wird bei Bedarf angelegt. \\
  \code{--summary <datei>} & Schreibt eine maschinenlesbare Übersicht aller Ergebnisse als JSON (Abschnitt~\ref{sec:summary}). \\
  \code{-j <n>}, \code{--threads <n>} & Höchstens $n$ Läufe einer Studie gleichzeitig (Standard: alle Kerne). \\
  \code{--progress} & Meldet während Transienten und Studien Zeilen \code{progress <Analyse> <Anteil>} auf stderr, höchstens eine je Prozent. \\
  \code{-h}, \code{--help} & Kurze Hilfe. \\
  \bottomrule
\end{tabularx}
\end{center}

Wo das Programm liegt:
\begin{itemize}
  \item nach dem Bauen: \datei{build/esim} (Windows:
    \texttt{build\textbackslash Release\textbackslash esim.exe}),
  \item im Programmpaket des Editors: unter Windows neben
    \datei{esim-editor.exe}, unter macOS in
    \datei{esim-editor.app/Contents/MacOS/esim}.
\end{itemize}

Wer es oft braucht, nimmt den Ordner in den Suchpfad (\code{PATH}) auf.

\section{Konsolenausgabe}

\begin{shell}
Common-emitter amplifier: bias point, frequency response, 1 kHz sine
11 components, 8 unknowns (6 node voltages, 2 branch currents, 0 equation variables)
[op] converged after 10 iterations (newton)
  v(b)   = 2.04845
  v(c)   = 5.52274
  v(e)   = 1.38503
  v(c,e) = 4.13771
  i(VCC) = -0.00158988
[ac] 1 .. 1e+08 Hz: 161 points -> results/bjt_amplifier.ac.csv
[noise] 10 .. 100000 Hz: 81 points -> results/bjt_amplifier.noise.csv
  output noise in the band: 2.13685e-05 rms
  referred to the input:    1.25945e-07 rms
    Q1.ic        98.71 %
    RC           0.7894 %
    ...
[tran] 0 .. 0.005 s, bdf (order up to 2): 506 steps (0 rejected, 0 Newton failures), ...
  thd = 0.0952244
  fourier, fundamental 1000 Hz -> results/bjt_amplifier.tran.fourier.csv
    v(out)  fundamental 1.66273, thd 9.52228 %
\end{shell}

Die erste Zeile ist der Titel, die zweite die Größe des Systems. Danach
folgt je Analyse eine Zeile mit ihren Kennzahlen und der Ergebnisdatei,
darunter Arbeitspunkt, Messgrößen, Rauschbeiträge oder Statistik.
Warnungen beginnen mit \code{warning:}, Hinweise mit \code{note:}. Ein Fehler
(\code{error: ...}) bricht ab; das Programm endet dann mit dem Rückgabewert
1, sonst mit 0. Skripte können das prüfen.

\section{Die Zusammenfassung}
\label{sec:summary}

\code{--summary ergebnis.json} schreibt alles, was ein anderes Programm über
den Lauf wissen muss:

\begin{json}
{
  "title": "RC-Tiefpass",
  "analyses": [
    { "type": "op", "name": "op", "strategy": "newton", "iterations": 1,
      "values": { "v(in)": 0.0, "v(out)": 0.0 } },
    { "type": "tran", "name": "tran", "file": ".../tiefpass.tran.csv",
      "measures": { "t_rise": 0.002145621433071054 },
      "stats": { "steps": 160, "rejected": 4, "newton_iterations": 240,
                 "newton_failures": 0, "factorizations": 155,
                 "max_order": 5, "events": 0 } },
    { "type": "ac", "name": "ac", "file": ".../tiefpass.ac.csv",
      "measures": { "fc": 159.18247552426536 } }
  ]
}
\end{json}

Je nach Analyse kommen hinzu: \code{sweep} (Parameter und Werte einer
Studie), \code{statistics} (Monte Carlo), \code{parameters} (Anpassung),
\code{normalized} und die absoluten Ableitungen (Empfindlichkeiten),
\code{fourier}. Der Editor und die Python-Schnittstelle lesen genau diese
Datei.

\section{Ergebnisse zeichnen: plot.py}

\datei{tools/plot.py} zeichnet die CSV-Dateien mit matplotlib
(\code{pip install matplotlib}):

\begin{shell}
python3 tools/plot.py results/buck_converter.tran.csv
python3 tools/plot.py results/bjt_amplifier.ac.csv -o bode.png
python3 tools/plot.py results/relaxation_oscillator.sweep.csv --cols "v\(c\)"
\end{shell}

Die Darstellung richtet sich nach der Kopfzeile der Datei: Zeitverläufe mit
getrennten Feldern für Spannungen, Ströme und übrige Größen, Frequenzgänge
als Betrag in dB und Phase über logarithmischer Frequenz, DC-Sweeps über der
ersten Spalte, Parameter-Sweeps mit einer Kurve je Parameterwert. Mehrere
Dateien werden übereinandergelegt.

\begin{center}
\begin{tabularx}{\textwidth}{@{}lL@{}}
  \toprule
  Option & Wirkung \\
  \midrule
  \code{-o bild.png} & Bild schreiben (PNG, PDF, SVG) statt ein Fenster zu öffnen \\
  \code{--cols <regex>} & nur Spalten, deren Name auf den regulären Ausdruck passt \\
  \code{--separate} & ein eigenes Feld je Größe \\
  \code{--group} & erste Spalte als Parameter eines Sweeps behandeln \\
  \code{--logx}, \code{--logy} & logarithmische Achsen \\
  \bottomrule
\end{tabularx}
\end{center}

\section{Skripte}

Weil Eingabe und Ausgabe Textdateien sind, lässt sich esim gut in Skripte
einbinden. Ein Beispiel für die Shell, das alle Beispiele rechnet:

\begin{shell}
for f in examples/*.json; do
  esim "$f" -o results || echo "FEHLER in $f"
done
\end{shell}

Für alles, was über einfache Schleifen hinausgeht – Netzlisten erzeugen,
Ergebnisse auswerten, Optimierungen –, ist die Python-Schnittstelle des
nächsten Kapitels bequemer.
